\section{Two independent half-permutations}\label{sec:two-halves}

The previous constructions extract contraction from repeated active intervals. This construction uses two long independent blocks. For the second block, reveal only which slots contain each of two colors at every layer. That observation leaves independent internal permutations of the two colors. We will trim their large tuple marginals and combine them with inverse-image information from the first block. The two orientations have different roles and must both be retained.

\subsection{The two precise inputs}

Here we use two companion results. We state their full hypotheses so that the conditional law to which each applies is explicit. Write $\operatorname{inj}(I,J)$ for the set of injections from a finite ordered set $I$ into a finite set $J$.

The conditional path theorem~\cite[Theorem~4.1]{conditional-information} has the following form. Fix $0<p<1$ and define
\begin{equation}\label{eq:tabloid-parameters}
 \rho_p=\frac{3+p}{4},\qquad c_p=-\frac18\log_2\rho_p,
 \qquad b=\lceil10/c_p\rceil.
\end{equation}
For every deterministic starting deck, every specified base set of $r_0$ labels, and every ordered list of $k$ additional labels satisfying $r_0+k-1\le pn$, after $2b$ sweeps,
\begin{equation}\label{eq:conditional-input}
 \E_{\mathcal H}\left\|
 \mathcal L(\text{list endpoints}\mid\mathcal H)
 -\operatorname{Unif}\operatorname{inj}([k],V\setminus E_{\mathcal H})
 \right\|_{\TV}\le k n^{-4}.
\end{equation}
Here $\mathcal H$ is the full labelled base-path history and $E_{\mathcal H}$ its endpoint set. The estimate holds for every cyclic coordinate order, including the reverse order. It also holds after coarsening $\mathcal H$, provided the endpoint available set remains determined. All thresholds depend only on the fixed $p$.

The abstract conditional-product transfer~\cite[Theorem ``Conditional-product transfer'']{partial-fourier} concerns $V=A\sqcup C$ with $|A|=|C|=h$, and $H=S_A\times S_C$. Partition each half into $L$ equal buffers, where $L\ge2C_0$ and $C_0=2000000$. Suppose a subprobability $\xi$ on $S_V$ has both ordered inverse-image marginals $x^{-1}|_A$ and $x^{-1}|_C$ pointwise at most twice uniform. For every environment $z$, let $\nu_A^z,\nu_C^z$ be subprobabilities on the two half-groups. Assume that, in each factor and for every buffer, the ordered image tuple on the complement of that buffer has marginal at most twice uniform. For each $z$, choose an arbitrary deterministic permutation $t_z\in S_V$. If $Z$ has any probability law, put
\[
 \omega=\E_Z\big[\delta_{t_Z}*(\nu_A^Z\otimes\nu_C^Z)\big],
 \qquad \theta=\omega*\xi.
\]
Then, for all sufficiently large $h$, every irreducible representation of $S_V$ of degree $D>1$ satisfies
\begin{equation}\label{eq:tabloid-transfer}
 \|\widehat\theta\|_{\op}\le\sqrt{22}\,D^{-1/40}.
\end{equation}
The Hilbert--Schmidt norms in that theorem use ordinary, unnormalized trace. For either $J\in\{A,C\}$, let $\nu_J$ be a factor satisfying the buffer caps and let $\gamma$ be an irreducible representation of $S_J$, of degree $D_\gamma$. The underlying buffer estimate is
$\|\widehat\nu_J(\gamma)\|_{\HS}^2/D_\gamma\le\min(1,2D_\gamma^{-1/2})$.
For irreducibles $\alpha$ of $S_A$, $\beta$ of $S_C$, and $\lambda$ of $S_V$, write $c_{\alpha\beta}^{\lambda}$ for the multiplicity of $\alpha\otimes\beta$ in the restriction of $\lambda$ to $H$. These multiplicities are retained through the estimate
$c_{\alpha\beta}^{\lambda}\le\min(D_\alpha,D_\beta)$, where each $D$ denotes the corresponding degree.
For the final Fourier sum we will use the local reciprocal-square estimate in Lemma~\ref{lem:degree-sum}.

These are separate inputs. The first supplies conditional tuple information from the shuffle; the second converts explicit caps and a conditional product into an operator bound. The product property will now be proved for the actual shuffle environment.

\subsection{Inverse-image caps for the first block}

\begin{theorem}[Two-half construction]\label{thm:two-halves}
Set $L=2^{23}$, $p=1-1/(2L)$, and choose $b$ by~\eqref{eq:tabloid-parameters}. For $n=2^d$, the full-deck law from any deterministic start approaches uniform in total variation after $400bd$ physical shuffles.
\end{theorem}

For these choices, $L$ divides $h=n/2$ for all sufficiently large $d$. Fix two halves $A,C$ of $V$ and partition each into $L$ buffers of size $h/L$. Let $X,W$ be independent permutations each generated by $2b$ forward sweeps, and let $\mu$ be the law of $WX$.

The inverse of $X$ is generated by the independent involutive layers in reverse order. Apply~\eqref{eq:conditional-input} with empty base and a list consisting of $A$, and then of $C$. Since $h-1\le pn$, each ordered inverse-image half-tuple has distance at most $h n^{-4}$ from uniform.

We use the following elementary truncation rule repeatedly. If probability laws $P,Q$ are on one finite set, then
\[
 P\{v:P(v)>2Q(v)\}\le2\|P-Q\|_{\TV},
\]
since $P(v)\le2(P(v)-Q(v))$ on those atoms and the total positive excess is the variation distance. Delete from the law of $X$ every permutation whose inverse-image tuple on either half is an atom of more than twice its uniform probability. Call the resulting subprobability $\xi$. It obeys both required pointwise caps, and
\begin{equation}\label{eq:x-mass}
 1-\xi(S_n)\le4h n^{-4}=2n^{-3}.
\end{equation}
The deletion for the other half cannot increase a marginal. Thus one common $\xi$ has the two caps simultaneously.

\subsection{The occupancy environment and the conditional product}

Color the positions at the start of $W$ by membership in $A$ or $C$, regarding those starting positions as distinct labels for this block. Let $Z$ record the set of positions occupied by each color after every layer, including the initial and final boundary. It records no labels within a color.

\begin{lemma}[Independent internal permutations]\label{lem:color-product}
For every feasible occupancy history $Z=z$, all mixed-color switch values are fixed and all same-color switch values are independent fair bits. Choose $t_z$ to send the increasing list of $A$ to the increasing list of its endpoint color set and likewise for $C$. Then
\[
 W=t_z(Y_A\times Y_C),
\]
where, conditional on $Z=z$, $Y_A\in S_A$ and $Y_C\in S_C$ are independent. Their conditional laws depend only on $z$.
\end{lemma}
\begin{proof}
At a mixed-color edge, the next pair of colors determines whether its coin is zero or one. At a same-color edge, either coin gives the same outgoing colors. Conversely, once the coins specified at the mixed edges of $z$ are fixed, induction on layers forces precisely the occupancy history $z$, independently of all same-color coins. Thus the event $Z=z$ is a cylinder in the time-edge array. Its unprescribed bits remain mutually independent.

A label of color $A$ uses fixed transports at mixed edges and only the independent bits on $A$--$A$ edges otherwise. The analogous assertion for color $C$ uses the disjoint array on $C$--$C$ edges. For fixed $z$ those two sets of time-edge coordinates are deterministic. Applying the slot-bijection construction separately to the two colors gives bijections onto their endpoint sets. The stated $t_z$ identifies these sets with the original halves, giving the independent permutations $Y_A,Y_C$. In particular no independence of the color history itself from the internal permutations is being used.
\end{proof}

Write $\mu_A^z,\mu_C^z$ for these two conditional probabilities. Fix a buffer $B\subset A$. Apply~\eqref{eq:conditional-input} to $W$ with all $C$ labels as base and the list $A\setminus B$ as the additional labels. Here
\[
 r_0+k-1=h+(h-h/L)-1
        =n-\frac n{2L}-1\le pn.
\]
The error is at most $(h-h/L)n^{-4}$. The full labelled $C$ paths determine the color occupancy history $Z$, and $Z$ still determines their endpoint set. Hence the allowed coarsening in~\eqref{eq:conditional-input} gives that same expected error conditional on $Z$. Applying the deterministic bijection $t_Z^{-1}$ sends the comparison law to uniform ordered distinct images in $A$. We conclude
\begin{equation}\label{eq:color-buffer-error}
 \E_Z\left\| (\mu_A^Z)_{A\setminus B}
          -(U_A)_{A\setminus B}\right\|_{\TV}
 \le(h-h/L)n^{-4}.
\end{equation}
The same argument applies to every buffer of $C$, interchanging the colors. The endpoints' available set is essential here: ordinary unconditional list mixing would not imply~\eqref{eq:color-buffer-error}.

For each $z$, delete from $\mu_J^z$ all permutations whose ordered image tuple outside any buffer exceeds twice its uniform probability, for $J=A,C$. Denote the subprobabilities by $\nu_J^z$ and their masses by $a_J(z)$. The truncation rule and~\eqref{eq:color-buffer-error} give
\[
 \E_Z\big[(1-a_A(Z))+(1-a_C(Z))\big]
 \le4L(h-h/L)n^{-4}=(2L-2)n^{-3}.
\]
Every surviving factor satisfies each cap for its particular $z$, even if the loss at that environment is large. Lemma~\ref{lem:color-product} now permits the product of these two deletions: define
\[
 \omega=\E_Z[\delta_{t_Z}*(\nu_A^Z\otimes\nu_C^Z)].
\]
It is dominated pointwise by the true law of $W$. Its missing mass is
$\E[1-a_Aa_C]$, bounded by the preceding display because
$1-a_Aa_C\le(1-a_A)+(1-a_C)$.

We have now established every conditional-law hypothesis of the transfer: the environment has its original probability law, each conditional measure is an actual product, each factor has all its buffer-complement caps, and the first block has the two inverse-image caps. Define $\theta=\omega*\xi$. Independence of the original blocks and positivity of convolution give $\theta\le\mu$ pointwise, while
\begin{equation}\label{eq:tabloid-mass}
 1-\theta(S_n)\le2L n^{-3}=o(1).
\end{equation}
The integer $L$ is large but fixed before $d$ tends to infinity.

\subsection{Fourier completion and physical time}

Applying~\eqref{eq:tabloid-transfer} gives the degree-power operator bound for every $D_\lambda>1$. We include the final deduction to keep the mass, parity and time conventions visible in this application.

The original block law $\mu$ has expected sign zero. Indeed, condition on all switch coins except any one coin in a physical layer. Flipping that fair coin composes the final permutation with a conjugate transposition and reverses its sign. From $\theta\le\mu$ and~\eqref{eq:tabloid-mass},
\[
 |\widehat\theta(\sgn)|\le1-\theta(S_n)\le2L n^{-3}.
\]
Let $j=100$ and $q=\theta(S_n)$. The trivial Fourier coefficient of
$\theta^{*j}-U_n$ is $q^j-1=o(1)$ and its sign coefficient is
$\widehat\theta(\sgn)^j=o(1)$. For all other irreducibles, without assuming their matrices commute or are normal,
\[
 \sum_{D_\lambda>1}D_\lambda
       \|\widehat\theta(\lambda)^j\|_{\HS}^2
 \le22^j\sum_{D_\lambda>1}D_\lambda^{2-2j/40}
 =22^{100}\sum_{D_\lambda>1}D_\lambda^{-3}=o(1).
\]
We used submultiplicativity of the operator norm,
$\|B\|_{\HS}^2\le D\|B\|_{\op}^2$, and Lemma~\ref{lem:degree-sum}, since $D^{-3}\le D^{-2}$.
Plancherel and Cauchy--Schwarz imply
$\|\theta^{*100}-U_n\|_1=o(1)$. Restoring the missing measure costs
$1-q^{100}=o(1)$, so $\|\mu^{*100}-U_n\|_{\TV}\to0$.

Each of $X,W$ costs $2bd$ physical shuffles. One macro-step therefore costs $4bd$, and the $100$ independent macro-steps cost $400bd$. They concatenate at sweep boundaries and hence have exactly the physical shuffle law at that time. Translation by any fixed initial deck preserves total variation. This proves Theorem~\ref{thm:two-halves}.
